\section{From a Learned Continuation to Economic Accuracy}\label{sec:r19decision}
\subsection{Centered approximation error and the selected decision}
Fix an anchor continuation $S_0$ and its approximation $\widehat S_0$.
For each task, let $e=\widehat S_0-S_0$ and
$\Delta_\theta=S_\theta-S_0$. The implemented action $\widehat a$ satisfies
\begin{equation}
 r_\theta(s,\widehat a)+\widehat S_0(T_\theta(s,\widehat a))
 \geq \sup_{a\in A_\theta(s)}
 \{r_\theta(s,a)+\widehat S_0(T_\theta(s,a))\}-\delta_\theta(s).
 \label{eq:r19greedy}
\end{equation}
Here $\delta_\theta(s)\geq0$ is an actual optimization allowance, not a
training-loss tolerance. For a function on the feasible actions, write
$\operatorname{osc}_A f=\sup_A f-\inf_A f$.

\begin{proposition}[Continuation error, transport, and decision loss]
\label{prop:r19decision}
Suppose the suprema in \eqref{eq:r19objective} and
\eqref{eq:r19greedy} are finite. Then
\begin{equation}
 L_\theta(s,\widehat a)\leq\delta_\theta(s)
  +\operatorname{osc}_{A_\theta(s)}(e\circ T_\theta)
  +\operatorname{osc}_{A_\theta(s)}(\Delta_\theta\circ T_\theta).
 \label{eq:r19osc}
\end{equation}
For a finite action set $\{a_1,\ldots,a_m\}$ with $m\geq2$ and fixed
positive weights $\pi_j$ summing to one, define
\begin{align*}
 R_\pi(s)&=\sum_j\pi_j(e(T_\theta(s,a_j))-\bar e_\pi(s))^2,\\
 \bar e_\pi(s)&=\sum_j\pi_j e(T_\theta(s,a_j)),
 &C_\pi&=\max_{i\ne j}(\pi_i^{-1}+\pi_j^{-1}).
\end{align*}
If $|\Delta_\theta|\leq b_\theta$ on the queried postdecision set, then
\begin{equation}
 L_\theta(s,\widehat a)\leq\delta_\theta(s)
                       +\sqrt{C_\pi R_\pi(s)}+2b_\theta.
 \label{eq:r19riskloss}
\end{equation}
Both inequalities hold for each fitted object and each realized cache.
Expectations over training or cache randomness may subsequently be taken.
\end{proposition}

The proof is in Section~\ref{app:r19proofs}. An additive continuation error
common to all feasible actions has no effect on the decision. This is why
an absolute squared prediction error may be large even when decisions are
accurate. Conversely, the risk at only the NBO/reference action pair cannot
control the losses of another method choosing elsewhere. Our new diagnostic
uses a common exogenous action set and, separately, the arguments selected
by each predictor. Inequality~\eqref{eq:r19riskloss} is a bound for the common
finite set; it is not applied to the entire scalar interval using just three
sampled actions.

If the true finite-menu best action has gap $g>0$ from every competitor,
then a bound strictly smaller than $g$ implies that the implemented action
is best. Without such a margin, arbitrarily small centered prediction error
can change the identity of the selected action. Neither statement orders two
methods' realized payoffs from their risks alone.

\subsection{A nontrivial reuse region when the future changes}
Let $X_{k+1}^\theta=F_{\theta,k}(X_k^\theta,Z_k)$ describe the future
under the policy being evaluated. A common-innovation coupling need not
make the transition maps identical. On an invariant state domain, assume
\begin{align}
 \|F_{0,k}(x,z)-F_{0,k}(x',z)\|&\leq L_k\|x-x'\|,
 &\|F_{\theta,k}(x,z)-F_{0,k}(x,z)\|&\leq\epsilon_{\theta,k},\notag\\
 |g_{0,k}(x)-g_{0,k}(x')|&\leq H_k\|x-x'\|,
 &|g_{\theta,k}(x)-g_{0,k}(x)|&\leq\eta_{\theta,k}.
 \label{eq:r19transportassumptions}
\end{align}
These bounds concern future rewards and transitions, not just the current
task. Include the terminal payoff as $g_{\theta,K}$ and use nonnegative
reward weights $\omega_k$. The initial postdecision state is shared.

\begin{theorem}[Continuation transport and refresh]
\label{thm:r19transport}
Set $d_0=0$ and
$d_{k+1}=L_kd_k+\epsilon_{\theta,k}$. Under
\eqref{eq:r19transportassumptions}, uniformly on the common postdecision
domain,
\begin{equation}
 |S_\theta-S_0|\leq
 b_\theta:=\sum_{k=0}^{K}\omega_k
                   (\eta_{\theta,k}+H_kd_k).
 \label{eq:r19transport}
\end{equation}
Thus the anchor may be reused at a task whenever the right side of
\eqref{eq:r19osc}, or the applicable finite-menu bound
\eqref{eq:r19riskloss}, is at most the desired economic tolerance.

More generally, suppose a recursive evaluation on a common invariant utility
domain has maps $V_{\theta,k}=\mathcal G_{\theta,k}(V_{\theta,k+1})$.
If these maps are sup-norm Lipschitz with constants $\beta_k$ on that domain,
$\|\mathcal G_{\theta,k}(v)-\mathcal G_{0,k}(v)\|_\infty\leq\nu_{\theta,k}$,
and $\|V_{\theta,K}-V_{0,K}\|_\infty\leq b_K$, then
\begin{equation}
 \|V_{\theta,0}-V_{0,0}\|_\infty\leq
 \sum_{k=0}^{K-1}\left(\prod_{j=0}^{k-1}\beta_j\right)\nu_{\theta,k}
       +\left(\prod_{j=0}^{K-1}\beta_j\right)b_K.
 \label{eq:r19recursive}
\end{equation}
Empty products equal one.
\end{theorem}

The theorem supplies a class of tasks with genuinely different futures.
For example, uniformly bounded transition and reward perturbations of an
anchor economy satisfy it on any verified invariant compact domain with
bounded Lipschitz constants. A single learned continuation then supports a
whole neighborhood of counterfactuals until its economic error allowance is
spent. The neighborhood is conditional on the displayed stability bounds;
it is not a distribution-free assertion about all future policies. When
future rewards, constraints, or opponents' policies violate those bounds,
the evaluated continuation must be recomputed or its transport allowance
reestablished. For recursive utility the common utility domain is essential:
a local Lipschitz constant cannot be extended across a singular utility
boundary. For games, this comparison evaluates a specified deviation problem;
it does not replace the fixed-point or best-response conditions in the
Economic Applications companion.

\subsection{An explicit amortization class}
The preceding inequalities give a sufficient class in which a shared scalar
continuation supports a cheaper declared accuracy procedure, not just a
post hoc interpretation of successful cells. Fix a finite exogenous query
family with $K$ distinct postdecision arguments and a common continuation
$S(x)=\phi(x)^{\mathsf T}\beta$. The feature map $\phi:\mathcal X\to\R^p$
may be a frozen neural dictionary or a conventional basis. Its construction
cost is charged, and it is fixed before the fitting noises below. At design
points $x_i$, observe $Y_i=S(x_i)+\xi_i$, where the noises are conditionally
independent, mean zero, and $\sigma$-sub-Gaussian. Let $\Phi$ have rows
$\phi(x_i)^{\mathsf T}$ and assume
\[
 G_n=\Phi^{\mathsf T}\Phi\text{ is invertible},\qquad
 \max_{x\text{ queried}}\phi(x)^{\mathsf T}G_n^{-1}\phi(x)\leq\ell/n.
\]
For example, $G_n/n\succeq cI$ and $\|\phi(x)\|^2\leq B_\phi^2$
imply $\ell=B_\phi^2/c$. These are checkable design conditions, not a
claim about an arbitrary trained representation.

\begin{corollary}[A sufficient learned-continuation work advantage]
\label{cor:r19amortization}
Fix $\varepsilon>\delta\geq0$ and $0<\alpha<1$. Least-squares continuation
fitting and $\delta$-accurate finite-menu maximization have loss at most
$\varepsilon$ on every query with probability at least $1-\alpha/2$ if
\begin{equation}
 n\geq \frac{8\sigma^2\ell}{(\varepsilon-\delta)^2}
                    \log\frac{4K}{\alpha},
 \label{eq:r19learnbudget}
\end{equation}
provided the displayed full-rank and leverage conditions hold. Direct
sample-average predictions of the same continuation have the same guarantee
when they use at least
$t=\lceil8\sigma^2\log(4K/\alpha)/(\varepsilon-\delta)^2\rceil$
observations per distinct argument, under the analogous scalar noise bound.
Common innovations across arguments are allowed.

Suppose the complete declared learned procedure costs
$F(n)+Kc_N+C(q)$ and the declared direct procedure costs
$C_0(t)+Kt c_R+C(q)$, with all fitting, caching, selection and verification
included. Both procedures meet the target jointly with probability at least
$1-\alpha$, and the learned procedure is strictly cheaper whenever
\begin{equation}
 K(t c_R-c_N)>F(n)-C_0(t).
 \label{eq:r19explicitadvantage}
\end{equation}
Repeated identical arguments are counted once in $K$. If verification or
selection costs differ, their actual difference enters the right side.
\end{corollary}

The class contains nonlinear continuations represented by finite smooth
neural features, not only linear economic models. The leverage bound keeps
fitting accuracy uniform as queries are reused. A bounded representation
bias or transported future adds its explicit allowance to the decision bound
and reduces the available margin in \eqref{eq:r19learnbudget}. The result
compares two fully specified sufficient-accuracy procedures; it is not a
lower bound on the best possible simulation method. A Raw procedure may
itself fit a surrogate, and then belongs to the learned comparison class.
In particular, this corollary neither certifies the assumptions for the
experiment's jointly trained network nor asserts that its conventional
surrogates must be slower. It provides an analytical reuse class that the
numerical results need not be asked to establish by extrapolation.

\subsection{Actual stopping and the amortization threshold}
Consider a feasible scalar action in $[\ell,u]$ and a twice differentiable
true objective with $Q''\leq-\kappa<0$. Define the directional residual
\[
 r(a)=\begin{cases}
 \max\{Q'(a),0\},&a=\ell,\\
 \max\{-Q'(a),0\},&a=u,\\
 |Q'(a)|,&\ell<a<u.
 \end{cases}
\]
An interval enclosure of $Q'(a)$ and a lower bound on $\kappa$ give a
computable upper bound $U(a)$ in the following statement.
\begin{proposition}[Certified stopping and complete work]\label{prop:r19stopping}
For every feasible scalar action,
\begin{equation}
 \sup_{b\in[\ell,u]}Q(b)-Q(a)\leq\frac{r(a)^2}{2\kappa}.
 \label{eq:r19residual}
\end{equation}
For a finite sequence of candidate-fitting stages, let a service stop at the
first stage whose mean outward bound is at most $\varepsilon$. If each bound
is a deterministic enclosure for the true task objectives, the returned
service satisfies the target regardless of the preceding failed checks.
If checks instead have conditional error allocations $\alpha_j$ given the
past and $\sum_j\alpha_j\leq\alpha$, the analogous conclusion holds with
probability at least $1-\alpha$.

The realized cost of a service stopping at stage $J$ is the sum of all
initialization, fitting or cache construction, querying, verification, and
durable-record costs through $J$. In an exact fixed-plus-linear cost model
$W_m(q)=F_m+q c_m$, where $c_m$ includes marginal verification and all
attempted stages, two methods that both meet the same target satisfy
\begin{equation}
 W_N(q)<W_R(q)\quad\Longleftrightarrow\quad
 q(c_R-c_N)>F_N-F_R.
 \label{eq:r19break}
\end{equation}
When $c_R>c_N$, this gives the strict break-even threshold. If only cost
intervals are available, a guaranteed ranking requires an upper bound for
one method below a lower bound for the other, not two upper bounds.
\end{proposition}

The cost model is an interpretation, not an imposed shape for the measured
frontier. Fitting stages may change with query volume, and verification can
dominate marginal prediction work. We therefore execute every volume rather
than extrapolate a line from a single 384-query comparison. The new experiment
runs each method's frozen stage sequence; it does not search globally for a
winner and then claim to have paid only that winner's bill. All comparison
services and shared startup are reported separately from the service that
an economist would deploy after choosing a method.
